$b_n = a_{2n}/a_n$ とおく. 
%$\alpha = 2 + \sqrt{5}$, $\beta=2-\sqrt{5}$とおくとき, $a_n = \frac{1}{2\sqrt{5}}(\alpha^n - \beta^n)$である(帰納法から分かる). よって
$b_n = \frac{\alpha^{2n} - \beta^{2n}}{\alpha^n - \beta^n} = \alpha^n + \beta^n$ である. 便宜上, $b_0 = 2$とおくと式と矛盾しない. $b_n$は漸化式$b_{n+2} = 4b_{n+1} + b_n$ ($n=0,1,\dots$), $b_0 = 2$, $b_1 = 4$を満たす. よって帰納的に$b_n$は偶数であり, とくに$b_n$が整数であることが従う. 漸化式を$\bmod{8}$することで $b_{n+2} \equiv b_n \pmod{8}$が従う($4b_{n+1}$は8の倍数であることに注意). よって初期値を見ることで, $n$が偶数ならば$b_n \equiv 2\pmod{8}$であり, $n$が奇数ならば$b_n\equiv 4\pmod{8}$であるので, 2で割り切れる回数は
\begin{center}
    \emph{$n$が奇数であるとき $2$回, \qquad $n$が偶数であるとき $1$回.}
\end{center}
$a_n$は漸化式$a_{n+2} =4a_{n+1} + a_n$, $a_0 = 0, a_1 = 1$を満たす.漸化式から帰納的に整数であることに注意する. まず$n$が奇数の場合, 漸化式より$a_{n+2}\equiv a_n \equiv\dots \equiv a_1 =1 \pmod{4}$であるので, $a_n$は奇数である. 階乗数は1を除いて偶数であるから, この場合は$a_n=1$となる必要がある. 数列$a_n$は$n\geq 0$で増加し, $a_1 = 1$であるから, この場合$n=1$のみが得られる.

以降は$n$が偶数でかつ$m\geq 2$の場合を考えるが, $n$を1以上の整数$e$と奇数$k$によって$n=2^e k$と表せるとする. このとき,
\[
a_{n} = a_{2^e k} = b_{2^{e-1}k}a_{2^{e-1}k} = \dots = (b_{2^{e-1}k}b_{2^{e-2}k} \dots b_{2k})b_ka_k
\]
となる. ただし$e=1$のときは最右辺の括弧を1とみなす. 最右辺の2で割り切れる回数について調べると, 各$b_{2^{j}k}$ ($1\leq j\leq e-1$) は(1)から2で1回だけ割り切れる. $k$は奇数なので$b_k$は2で2回割り切れる. $a_k$は前半で議論したように, $k$が奇数なので$a_k$も奇数である. よって$a_n$が2で割り切れる回数はちょうど$e-1 + 2 = e+1$であるということになる.

$m!$が2で割り切れる回数は $\sum_{k=1}^{\infty} \parenc{\frac{m}{2^k}}$と表せる. よって$e + 1 = \sum_{k=1}^{\infty} \parenc{\frac{m}{2^k}} \geq \frac{m}{2}$が必要である(最後の不等号は$m\geq 2$で成り立つ). したがって,
\bealn{
(2e + 2)! &\geq m! &(m\leq 2e+2)\\
&= a_n & (\text{条件式}) \\
&\geq a_{2^e} & (\text{数列$a_n$の増加性})\\
&> \frac{1}{2\sqrt{5}}(\alpha^{2^e} - 1) & (0 < \beta^{2^e} < 1)\\
&> \frac{1}{8}(4^{2^e} - 1) & (4 < \alpha, \quad 2\sqrt{5} < 8)
}
が従う. これが成り立つ$e$の範囲について調べる. まず$e=1$では$4! > 15/8$となり成立する. 次に$e\geq 2$では,
\[
(2e+2)! > 4^{2^e - 1.5} - \frac{1}{8}
\]
と見て$4^{2^e - 1.5} = 2^{2^{e+1} - 3}$が整数であることに注意すると, 上の不等式が成り立つならば
\begin{equation} \label{eq:7-2-1}
    (2e+2)! \geq 2^{2^{e+1} - 3}
\end{equation}
も成り立つ. この不等式が$e\geq 4$で成り立たないことを数学的帰納法で証明する. まず$e=4$のときは$10! = 3628800 < 2^{29}$なので正しい($2^{22} = 4\cdot 1048576$と比較する). 次に, ある$e\geq 4$において$(2e + 2)! < 4^{2^e - 3}$が正しいと仮定するとき, 両辺に$2^{2^{e+1}}$を掛けて
\[
2^{2^{e+1}}\cdot (2e + 2)! < 2^{2^{e+2} - 3}
\]
が成立する. よって$(2e + 3)(2e + 4) < 2^{2^{e+1}}$を示せば$e+1$の場合も正しいことが分かる. この不等式は, $2^e \geq 1 + e + {}_{e}{\mathrm{C}}_{2} + e + 1 \geq 2e + 4 > 2e$に注意すれば, 次の評価から正しいとわかる:
\[
(2e + 3)(2e + 4) < 2^e\cdot 2^e = 2^{2e} < 2^{2^e} < 2^{2^{e+1}}.
\]
よって数学的帰納法により, $e\geq 4$では\cref{eq:7-2-1}は成り立たず, $e = 1,2,3$の場合を見ればよいことになる. ここで割り切れる回数に戻ると, $m!$が2で割れる回数が$e+1$だったので, $m!$は高々2で4回しか割れない. よって$m\leq 7$であり, $a_n \leq 5040$である. $a_n$の漸化式から
\[
a_0 = 0, a_1 = 1, a_2 = 4, a_3 = 17, a_4 = 72, a_5 = 305, a_6 = 1292, a_7 = 5473,\dots
\]
であり, 以降は増加するので$n\geq 7$において解は存在しない. また, 上の中で階乗数であるものは1のみである. よって, \emph{求める組は$(1,1)$のみであることが示された.}
